\section{Positive shifts of density matrices}
\label{sec:shifted_density_powers}

Let $I$ be a finite index set and write
$\mathcal D_I=\{x\in\C^{I\times I}:x\ge0,\ \tr(x)=1\}$.
All norms in this section are operator norms for the Euclidean norm on
$\C^I$. For $b>0$, real powers of $x+b\Id$ are taken on the full space.

% Independent finite-dimensional support for the regularized filters in
% OpenAI, Polynomial PEPS approximation of gapped square-grid ground states,
% September 24, 2026, 03-patches.tex, lines 51--114, at
% adc7f1241b42e322a6451854ab7e4b4c146bf78a.
% These entries do not assert the ordered-product objective, its minimum,
% stationarity, energy estimates, or Proposition 4.1.

\begin{theorem}[Compactness and existence of density matrices]
    \label{thm:density_finite_index_compact}
    \lean{Matrix.isCompact_setOf_posSemidef_trace_eq_one,
        Matrix.setOf_posSemidef_trace_eq_one_nonempty}
    \leanok
    The set $\mathcal D_I$ is compact. If $I$ is nonempty, then
    $\mathcal D_I$ is nonempty.
\end{theorem}

\begin{proof}
    \leanok
    \uses{thm:density_compact, thm:density_nonempty}
    Choose a bijection from $\{0,\ldots,|I|-1\}$ to $I$.
    Reindexing rows and columns preserves positive semidefiniteness and
    trace, and is continuous. Thus $\mathcal D_I$ is the continuous image
    of the compact set of density matrices in dimension $|I|$ from
    Theorem~\ref{thm:density_compact}. If $I$ is nonempty, then $|I|>0$,
    and reindexing a density matrix supplied by
    Theorem~\ref{thm:density_nonempty} proves existence.
\end{proof}

\begin{theorem}[Positive identity shift]
    \label{thm:positive_identity_shift}
    \lean{Matrix.PosSemidef.add_smul_one_posDef}
    \leanok
    If $x\ge0$ and $b>0$, then $x+b\Id$ is positive definite.
\end{theorem}

\begin{proof}
    \leanok
    The matrix $b\Id$ is positive definite. Adding the positive
    semidefinite matrix $x$ preserves positive definiteness.
\end{proof}

\begin{theorem}[Continuity of shifted real powers]
    \label{thm:shifted_real_power_continuous}
    \lean{Matrix.continuousOn_add_smul_one_rpow}
    \leanok
    For every fixed $b>0$ and $r\in\R$, the function
    $x\mapsto(x+b\Id)^r$ is continuous on the positive semidefinite cone.
\end{theorem}

\begin{proof}
    \leanok
    \uses{thm:positive_identity_shift}
    The affine map $x\mapsto x+b\Id$ is continuous and, by
    Theorem~\ref{thm:positive_identity_shift}, takes the positive
    semidefinite cone into the strictly positive matrices. Continuous
    functional calculus makes $y\mapsto y^r$ continuous on that set.
    Composing these two maps proves the assertion.
\end{proof}

\begin{theorem}[Spectrum of a shifted density matrix]
    \label{thm:shifted_density_spectrum}
    \lean{Matrix.PosSemidef.spectrum_add_smul_one_bounds}
    \leanok
    If $x\in\mathcal D_I$ and $b>0$, then
    $\spec(x+b\Id)\subseteq[b,1+b]$.
\end{theorem}

\begin{proof}
    \leanok
    \uses{thm:positive_identity_shift}
    Positivity and the trace condition give
    $\|x\|\le\tr(x)=1$, hence $0\le x\le\Id$.
    Adding $b\Id$ yields
    $b\Id\le x+b\Id\le(1+b)\Id$.
    The matrix $x+b\Id$ is Hermitian by
    Theorem~\ref{thm:positive_identity_shift}, so these order inequalities
    give the corresponding lower and upper bounds on its spectrum.
\end{proof}

\begin{theorem}[Norm bounds for shifted density powers]
    \label{thm:shifted_density_power_norm}
    \lean{Matrix.PosSemidef.l2_opNorm_add_smul_one_rpow_le_of_nonpos,
        Matrix.PosSemidef.l2_opNorm_add_smul_one_rpow_le_of_nonneg}
    \leanok
    Let $x\in\mathcal D_I$, $b>0$, and $r\in\R$.
    If $r\le0$, then $\|(x+b\Id)^r\|\le b^r$.
    If $r\ge0$, then $\|(x+b\Id)^r\|\le(1+b)^r$.
    These bounds include $r=0$ and singular density matrices $x$.
\end{theorem}

\begin{proof}
    \leanok
    \uses{thm:positive_identity_shift, thm:shifted_density_spectrum}
    The norm estimate for continuous functional calculus bounds
    $\|(x+b\Id)^r\|$ by any uniform bound on $|t^r|$ over
    $t\in\spec(x+b\Id)$. By
    Theorem~\ref{thm:shifted_density_spectrum}, $b\le t\le1+b$.
    The function $t\mapsto t^r$ is positive there, decreasing for
    $r\le0$ and increasing for $r\ge0$. Its respective bounds are
    therefore $b^r$ and $(1+b)^r$.
\end{proof}
